\subsection{Rischi organizzativi}

\subsubsection{RO1: Impegni universitari}
\label{rischio:RO1}
\begin{itemize}
    \item \textbf{Descrizione}: la preparazione per esami universitari può diminuire la disponibilità di uno o più membri del gruppo.
    \item \textbf{Pericolosità}: Media.
    \item \textbf{Probabilità}: Bassa.
    \item \textbf{Prevenzione}: è importante condividere le date in cui si verranno svolti gli esami e i periodi in cui si pensa di poter dare meno disponibilità per il progetto.
    \item \textbf{Mitigazione}: sarà compito del responsabile gestire il carico di lavoro e le attività da svolgere, basandosi sulle ore disponibili per ogni membro, in modo tale da garantire il rispetto delle scadenze. 
\end{itemize}

\subsubsection{RO2: Impegni personali}
\label{rischio:RO2}
\begin{itemize}
    \item \textbf{Descrizione}: un membro potrebbe garantire una minore disponibilità durante un periodo a causa di impegni di natura personale.
    \item \textbf{Pericolosità}: Bassa.
    \item \textbf{Probabilità}: Media.
    \item \textbf{Prevenzione}: è importante avvisare tempestivamente e con il maggior anticipo possibile i giorni in cui non si sarà disponibili. 
    \item \textbf{Mitigazione}: è compito del responsabile calcolare le ore disponibili del gruppo durante lo \glossario{sprint} e organizzare il carico di lavoro tenendo conto di tutti gli impegni precedentemente comunicati dai compagni. 
\end{itemize}

\subsubsection{RO3: Errori di pianificazione delle attività}
\label{rischio:RO3}
\begin{itemize}
    \item \textbf{Descrizione}: l'inesperienza dei membri del gruppo nel contesto della pianificazione di progetto è un fattore di rischio che può interferire sulla corretta realizzazione del progetto. 
    Nello specifico, la mancata esperienza nell'ambito può portare a una sottostima delle risorse necessarie allo sviluppo delle attività o a un errata valutazione delle tempistiche. 
    \item \textbf{Pericolosità}: Alta.
    \item \textbf{Probabilità}: Media.
    \item \textbf{Prevenzione}: continuo controllo del Piano di Progetto per monitorare le risorse utilizzate e il tempo in cui sono state impiegate.
    Aggiornamento tra i membri del gruppo tramite Telegram sul proseguo nello sviluppo delle attività.
    \item \textbf{Mitigazione}: in caso dovessero verificarsi difficoltà o ritardi è compito del responsabile rivedere il Piano di Progetto e modificare il numero delle risorse necessarie per la realizzazione delle attività durante lo \glossario{sprint} o, in casi gravi, la scadenza.
\end{itemize}

\subsubsection{RO4: Scarsa collaborazione di uno o più membri del gruppo}
\label{rischio:RO4}
\begin{itemize}
    \item \textbf{Descrizione}: rischio dovuto alla possibilità che uno o più membri impieghino scarso impegno nella realizzazione del progetto.
    \item \textbf{Pericolosità}: Alta.
    \item \textbf{Probabilità}: Bassa.
    \item \textbf{Prevenzione}: è stata da subito richiesta la massima trasparenza tra i membri del gruppo.
    Sarà compito del responsabile verificare che ogni membro abbia partecipato equamente nella realizzazione delle attività dello \glossario{sprint}.
    \item \textbf{Mitigazione}: Il responsabile dovrà contattare e avvisare il membro che, senza preavviso, non avrà impiegato le giuste risorse personali nello sviluppo delle attività dello \glossario{sprint}, in modo tale da chiarire eventuali complicanze o fraintendimenti. 
\end{itemize}

\subsubsection{RO5: Disaccordi all'interno del gruppo}
\label{rischio:RO5}
\begin{itemize}
    \item \textbf{Descrizione}: possono generarsi discussioni a causa di opinioni diverse all'interno del gruppo.
    \item \textbf{Pericolosità}: Media.
    \item \textbf{Probabilità}: Media.
    \item \textbf{Prevenzione}: importante che vi sia chiarezza e trasparenza durante le chiamate del gruppo, bisogna dichiarare da subito eventuali discordanze su idee e opinioni in modo da trovare un punto d'incontro più rapidamente possibile.
    \item \textbf{Mitigazione}: in caso in cui non si riesca a raggiungere un compromesso definitivo le varie opzioni su cui si sta discutendo verranno messe a votazione, verrà scelta l'opzione, o le opzioni, che hanno la maggioranza dei voti. 
\end{itemize}

\subsubsection{RO6: Problemi nello sviluppo della progettazione}
\label{rischio:RO6}
\begin{itemize}
    \item \textbf{Descrizione}: il gruppo potrebbe trovare difficoltà nello sviluppo della progettazione a causa dell'inesperienza dei membri o della complessità del progetto;
    \item \textbf{Pericolosità}: Alta.
    \item \textbf{Probabilità}: Alta.
    \item \textbf{Prevenzione}: è importante che i membri del gruppo effettuino allenamento e studio autonomo per migliorare le proprie conoscenze in ambito progettuale e che i dubbi che sorgono vengano chiariti il prima possibile.
    \item \textbf{Mitigazione}: in caso di difficoltà è importante che il gruppo si confronti e discuta le problematiche riscontrate, in modo da trovare una soluzione o in caso di non riuscirci, contattare il professore e/o il referente dell'azienda per ricevere supporto.
\end{itemize}

\subsubsection{RO7: Carico di lavoro squilibrato}
\label{rischio:RO7}
\begin{itemize}
    \item \textbf{Descrizione}: il carico di lavoro potrebbe non essere equamente distribuito tra i membri del gruppo, portando a un aumento del rischio di ritardi e malfunzionamenti;
    \item \textbf{Pericolosità}: Bassa.
    \item \textbf{Probabilità}: Media.
    \item \textbf{Prevenzione}: è importante che il responsabile verifichi che il carico di lavoro sia equamente distribuito tra i membri del gruppo e che ogni membro abbia la possibilità di esprimere le proprie opinioni riguardo le attività da svolgere;
    \item \textbf{Mitigazione}: in caso di squilibrio nel carico di lavoro il responsabile dovrà rivedere il \textit{Piano di Progetto} e riassegnare le attività in modo da garantire un carico di lavoro equo per tutti i membri del gruppo.
\end{itemize} 

\subsubsection{RO8: Difficoltà nel rispetto delle scadenze}
\label{rischio:RO8}
\begin{itemize}
    \item \textbf{Descrizione}: Un membro del gruppo potrebbe non rispettare le scadenze stabilite durante la pianificazione dello \glossario{sprint}, portando a ritardi nello sviluppo del progetto;
    \item \textbf{Pericolosità}: Media.
    \item \textbf{Probabilità}: Media.
    \item \textbf{Prevenzione}: È importante che il responsabile verifichi costantemente il rispetto delle scadenze, il proseguimento delle attività e che ogni membro del gruppo comunichi eventuali difficoltà riscontrate;
    \item \textbf{Mitigazione}: In caso di ritardi inevitabili il responsabile dovrà rivedere il \textit{Piano di Progetto} e riassegnare le attività in modo da garantire il rispetto delle scadenze.
\end{itemize}